\documentclass[a4paper]{article}

\usepackage{graphicx}
\usepackage{amsmath,amssymb}
\usepackage{minted}
\usepackage{multirow}
\usepackage{float}
\usepackage{todonotes}
\usepackage{forest}
\usepackage{xstring}
\usepackage{xcolor}
\usepackage[most]{tcolorbox}
\usepackage{xparse}

\usetikzlibrary{graphs,graphdrawing}
\usegdlibrary{layered}

% for external graphviz
\input{/home/donkarlo/Dropbox/repo/nd_ai_project/src/nd_ai/knoweldge_representation/language/formal/computer/programming/tex/latex/code_clips/external_graphviz.tex}

%for tags, comma separated \semtag{tag1, tag2}
\input{/home/donkarlo/Dropbox/repo/nd_ai_project/src/nd_ai/knoweldge_representation/language/formal/computer/programming/tex/latex/parts/control_sequence/kinds/semtag/semtag.tex}

% for links, just type \seehere
\input{/home/donkarlo/Dropbox/repo/nd_ai_project/src/nd_ai/knoweldge_representation/language/formal/computer/programming/tex/latex/code_clips/seehere.tex}

\title{}
\date{}

\begin{document}
    \maketitle
    
    
    \section{Evaluation Procedure for Temporal Discrepancy Localization}
        
        \subsection{Step 1: Train the forecasting model}
            Train the time-series forecasting model using data from the normal scenario.
        
        \subsection{Step 2: Estimate the error distribution on normal data}
            Run the trained model on the seen (normal) scenario and compute the distribution of prediction errors.
        
        \subsection{Step 3: Compute the discrepancy signal}
            For the evaluation scenario, compute the discrepancy between the prediction and the observed signal.
            
            \[
                e_t = |y_t - \hat{y}_t|
            \]
            
            If the signal contains multiple features:
            
            \[
                e_t = \frac{1}{F}\sum_{f=1}^{F} |y_{t,f} - \hat{y}_{t,f}|
            \]
            
            The resulting sequence
            
            \[
                e_1, e_2, \dots, e_T
            \]
            
            represents the discrepancy time series.
        
        \subsection{Step 4: Smooth the discrepancy signal}
            To reduce noise, apply a temporal smoothing window:
            
            \[
                \tilde{e}_t =
                \frac{1}{W}
                \sum_{i=t-W+1}^{t} e_i
            \]
        
        \subsection{Step 5: Define the expected unseen interval}
            Define the contiguous time interval where the scenario differs from the training data:
            
            \[
                I_{\text{novel}} = [t_a, t_b]
            \]
            
            The remaining time steps form the normal interval:
            
            \[
                I_{\text{normal}} =
                [1,T] \setminus I_{\text{novel}}
            \]
        
        \subsection{Step 6: Compute the mean discrepancy in each interval}
            
            \[
                \mu_{\text{novel}} =
                \frac{1}{|I_{\text{novel}}|}
                \sum_{t \in I_{\text{novel}}} e_t
            \]
            
            \[
                \mu_{\text{normal}} =
                \frac{1}{|I_{\text{normal}}|}
                \sum_{t \in I_{\text{normal}}} e_t
            \]
        
        \subsection{Step 7: Compute the discrepancy ratio}
            
            \[
                R =
                \frac{\mu_{\text{novel}}}{\mu_{\text{normal}}}
            \]
            
            If \(R\) is significantly larger than \(1\), the model error is concentrated within the unseen interval.
        
        \subsection{Step 8: Identify the temporal location of the largest discrepancy}
            
            \[
                t_{\max} =
                \arg\max_{t} e_t
            \]
            
            Ideally, \(t_{\max}\) should lie inside or near \(I_{\text{novel}}\).
        
        \subsection{Step 9: Define an error threshold from the training scenario}
            
            \[
                \theta =
                \mu_{\text{train error}}
                +
                3\sigma_{\text{train error}}
            \]
            
            For each time step:
            
            \[
                e_t > \theta
            \]
            
            indicates a high discrepancy.
        
        \subsection{Step 10: Evaluate temporal localization performance}
            
            Compute the fraction of time steps above the threshold inside and outside the unseen interval:
            
            \[
                recall_{\text{interval}} =
                \frac{
                    \#\{t \in I_{\text{novel}} : e_t > \theta\}
                }{
                    |I_{\text{novel}}|
                }
            \]
            
            \[
                false\ positive\ rate =
                \frac{
                    \#\{t \in I_{\text{normal}} : e_t > \theta\}
                }{
                    |I_{\text{normal}}|
                }
            \]
            
            A successful forecasting model should produce high discrepancy within the unseen interval while maintaining low discrepancy elsewhere.
    % \bibliographystyle{plain}
    % \bibliography{/home/donkarlo/Dropbox/repo/nd_shared_research/bibliography/bibliography.bib}
\end{document}